%! TEX root = main.tex

\begin{solution}{ex:em-1}
  (a)

  Using Coulomb's law, we have

  \[\abs{F}=k \frac{\abs{q_1q_2}}{r^2}\]
  The force on $q_3$ due to $q_1$:

  \[\abs{F_{31}}=k \frac{10.0 \cdot 10^{-12} C}{0.15m} \approx 0.6N\]
  The direction is away from $q_1$, since they are both positively charged. The force on $q_3$ due to $q_2$:

  \[\abs{F_{32}}=k \frac{15.0 \cdot 10^{-12} C}{0.25cm} \approx 0.54 N \]
  This time the force points towards $q_2$ since they are oppositely charged.
  (b)

  The net force is just the sum of the previous 2 forces in this case.
  \[\overrightarrow{F_{net}}=\overrightarrow{F_{31}}+\overrightarrow{F_{32}}\]
  It is clear that $\overrightarrow{F_{32}}$ and $\overrightarrow{F_{31}}$ both point in the positive x-axis direction. Thus:
  \[
  \begin{gathered}
        F_{net}=F_{31}+F_{32}\\
        F_{net}=0.6N+0.54N=1.14N
  \end{gathered}
\]
  (c)
  The electric field is given by:
  \[
    E=\frac{F_{net}}{q_0}
  \]
  We have the following equation:
  \[
    \begin{aligned}
      E&=0\\
      \frac{F_{net}}{q_0}&=0\\
      \frac{\frac{q_1q_0}{x^2}+\frac{q_2q_0}{(40-x)^2}}{q_0}&=0\\
      \frac{q_1}{x^2}+\frac{q_2}{(40-x)^2}&=0\\
      \frac{(40-x)^2q_1}{x^2}+q_2&=0\\
      \frac{(40-x)^2}{x^2}&=-\frac{q_2}{q_1}\\
      \frac{40-x}{x}&=\sqrt{-\frac{q_2}{q_1}}\\
      40-x&=x\sqrt{-\frac{q_2}{q_1}}\\
      40&=x(1+\sqrt{-\frac{q_2}{q_1}})\\
      x&=\frac{40}{1+\sqrt{-\frac{q_2}{q_1}}}\\
      x&\approx18cm
    \end{aligned}
  \]
\end{solution}
\begin{solution}{ex:em-2}
  The electric field strength is given by:
  \[
    E=\frac{U}{d}
    E=133 333 V/m
  \]
  (b) The force is given by
  \[
    F=qE
  \]
  So the acceleration of the electron on is just:
  \[
    \begin{gathered}
      a=\frac{eE}{m_e}\\
      a\approx2.3 \cdot 10^{16} \frac{m}{s^2}
    \end{gathered}
  \]
  (c) The electric potential between the 2 plates is given by:
  \[
    \Delta W_pot = q\Delta V
  \]
  We also have:
  \[
    \begin{gathered}
      U_{AB} = \frac{W_{AB}}{e}\\
      400V \cdot e = W_{AB}\\
      400V \cdot e = \frac{1}{2}m_e v^2\\
      v=\sqrt{\frac{2 \cdot 400V \cdot e}{m_e}}\\
      v\approx 11.8 \cdot 10^{12} m/s
    \end{gathered}
  \]
\end{solution}
\begin{solution}{ex:em-3}
  (a) Its capacitance is given by:
  \[
    C=\varepsilon_0\varepsilon_r \frac{A}{d}
  \]
  We'll use $\varepsilon_0 = 8.85 \cdot 10^{-12} F/m$ $\varepsilon_r = 4.0 F/m$, $A=0.005 m^2$ and $d=0.005 m$. We have:
  \[
    \begin{gathered}
      C\approx3.5 \cdot 10^{-11} F
    \end{gathered}
  \]
  (b)
  We have:
  \[
    Q=CU
  \]
  So using the capacitance obtained earlier we have:
  \[
    Q\approx8.5 \cdot 10^{-10}
  \]
  (c) The battery is disconnected, let's treat Q as fixed and U as variable. The capacitance is half of what it used to be, since the separation is doubled. We have:
  \[
    \begin{gathered}
      C_f=1.75 \cdot 10^{-11} F = \frac{C}{2}\\
      U_f=\frac{Q}{C_f}=48 V
    \end{gathered}
  \]
  The stored energy is given by:
  \[
    W_f=\frac{1}{2}C_fU_f^2\approx2.03 \cdot 10^{-8} J
  \]
  It is also clear that this energy is double of what it used to be. This is because separating the plates is doing work on the capacitor.
\end{solution}
\begin{solution}{ex:em-4}
  The total resistance is given by:
  \[
    \begin{gathered}
      R_t=R_1+R_{23}\\
      R_t=R_1+\frac{R_2R_3}{R_2+R_3}
      R_t=282 \Omega
    \end{gathered}
  \]
\end{solution}
